\documentclass[12pt]{article}
\usepackage[utf8]{inputenc}
\usepackage[T1]{fontenc}
\usepackage[margin=1in]{geometry}
\usepackage{setspace}
\usepackage{xcolor}
\usepackage[colorlinks=true, urlcolor=blue!60!black, linkcolor=blue!60!black]{hyperref}
\setlength{\parindent}{2em}
\setlength{\parskip}{1.25em}
\renewcommand{\baselinestretch}{1.0}
\begin{document}
\begin{center}
{\LARGE \textit{New York Times Co. v. United States}, 403 U.S. 713 (1971)\par}
\vspace{0.8em}
{\large Case Brief --- Four-Page Edition\par}
\end{center}
\vspace{1em}

\section{Memory Jogger}

In the Pentagon Papers Case, the Nixon Administration sought to stop the \textit{New York Times} and \textit{Washington Post} from publishing a classified Vietnam War study; a 6--3 \textit{per curiam} held that the Government had not carried the ``heavy burden'' required to justify a prior restraint, with all nine Justices writing separately (\href{https://supreme.justia.com/cases/federal/us/403/713/}{Justia}).

\section{Detailed Case Facts}

In 1967 Secretary of Defense Robert McNamara commissioned a classified, 47-volume history of U.S. decision-making in Vietnam through 1968, popularly known as the ``Pentagon Papers.'' In 1971 former defense analyst Daniel Ellsberg leaked copies to the \textit{New York Times} and later to the \textit{Washington Post}. On June 13, 1971, after months of internal review, the \textit{Times} began publishing excerpts and analysis. The Administration sued to halt further publication, asserting that disclosure would cause grave and irreparable injury to national security and foreign relations.

\section{Procedural History}

In the \textit{Times} case, Judge Murray Gurfein issued a temporary restraining order but then denied a permanent injunction; the en banc Second Circuit reversed, continued the restraint, and remanded for hearings on specific items. In the \textit{Post} case, Judge Gerhard Gesell denied an injunction and the en banc D.C. Circuit affirmed. Both cases reached the Supreme Court on June 24; it granted certiorari, entered stays on June 25, heard argument on June 26, and ruled on June 30, 1971.

\section{Judicial Votes}

\textbf{6--3}: the \textit{Times} judgment reversed, the \textit{Post} judgment affirmed, and all stays vacated.
\begin{itemize}
\item \textbf{Per curiam:} a brief unsigned opinion; no rationale commanded a majority.
\item \textbf{Concurring:} Justices Black (joined by Douglas), Douglas (joined by Black), Brennan, Stewart (joined by White), White (joined by Stewart), and Marshall.
\item \textbf{Dissenting:} Chief Justice Burger; Justice Harlan (joined by Burger and Blackmun); and Justice Blackmun.
\end{itemize}

\section{Holding}

Any system of prior restraints bears a ``heavy presumption against its constitutional validity,'' and the Government ``carries a heavy burden of showing justification for the imposition of such a restraint.'' The Government did not carry that burden, so the injunctions could not stand. The separate opinions supply the reasoning.

\section{Analysis of Opinions}

\textbf{Per curiam.} Relying on \textit{Bantam Books} and \textit{Near v. Minnesota}, the Court agreed with the lower courts that had refused injunctions.

\textbf{Black and Douglas.} ``No law'' means no prior restraint; each moment the injunctions remained was a ``flagrant, indefensible, and continuing violation of the First Amendment.'' Douglas added that Congress deliberately omitted ``publish'' from the Espionage Act's \S~793(e) and had rejected a 1917 censorship proposal.

\textbf{Brennan.} Restraints cannot rest on ``surmise or conjecture''; only proof that publication will ``inevitably, directly, and immediately'' cause a catastrophe like imperiling a transport at sea could suffice.

\textbf{Stewart.} Because the Executive's power over defense and foreign affairs is largely unchecked, a free press is the main democratic check. The Executive must keep its own secrets---``when everything is classified, then nothing is classified''---but courts may not enjoin publication absent direct, immediate, and irreparable damage.

\textbf{White and Marshall.} White conceded likely ``substantial damage to public interests'' but found no congressional authorization for an injunction, warning that criminal prosecution under the Espionage Act remained possible. Marshall, invoking \textit{Youngstown}, said courts may not ``make law'' Congress twice refused to enact.

\textbf{Dissents.} Burger protested the ``unseemly haste'' on an incomplete record. Harlan would confine review to whether the subject lay within the foreign-affairs power and a cabinet officer had personally found irreparable harm. Blackmun favored balancing and warned of the ``death of soldiers.''

\section{Examples: Future Applications}

\begin{enumerate}
\item \textbf{Same-Side (Barred): Memo on Exaggerated Intelligence.} A leaked historical memo showing officials exaggerated intelligence offers no concrete showing of grave and imminent harm.
\item \textbf{Same-Side (Barred): State Evacuation Plan.} A Governor's fear that publishing an inadequate hurricane plan could cause panic is ``surmise or conjecture,'' far short of Stewart's standard.
\item \textbf{Opposite-Side (Stronger Case): Live Convoy Coordinates.} Publishing live naval coordinates and cryptographic keys resembles Brennan's troop-transport example, though the opinions yield no reliable vote count, and prosecution under \S~798 is a separate question.
\item \textbf{Opposite-Side (Permitted): A Former Analyst's Memoir.} Under \textit{Snepp v. United States}, an employee who agreed to prepublication review may be enjoined and stripped of profits; she is not an outside newspaper.
\item \textbf{Fence-Sitter: A Narrow Injunction Statute.} A statute authorizing injunctions against naming covert officers abroad answers White's and Marshall's objection, but Black, Douglas, and Stewart's standard leave its validity unresolved.
\end{enumerate}

\section{Critique}

\textbf{Fractured guidance.} With nine opinions and no majority rationale, the case left open whether Congress could enact a narrow prior-restraint statute for military secrets, or whether a temporary restraint for \textit{in camera} review is ever permissible.

\textbf{Restraint versus punishment.} Critics note that White's distinction can be inverted in practice: the threat of a ten-year Espionage Act conviction may chill journalism more than an injunction proceeding in which a publisher can test legality before risking prison.

\section{Key Quotations}

\begin{quote}
``I am convinced that the Executive is correct with respect to some of the documents involved. But I cannot say that disclosure of any of them will surely result in direct, immediate, and irreparable damage to our Nation or its people. That being so, there can under the First Amendment be but one judicial resolution of the issues before us.''
\end{quote}
\noindent --- Justice Stewart, concurring, 403 U.S. at 730.

\begin{quote}
``The press was to serve the governed, not the governors. The Government's power to censor the press was abolished so that the press would remain forever free to censure the Government. The press was protected so that it could bare the secrets of government and inform the people.''
\end{quote}
\noindent --- Justice Black, concurring, 403 U.S. at 717.

\begin{quote}
``Great cases, like hard cases, make bad law. ... With all respect, I consider that the Court has been almost irresponsibly feverish in dealing with these cases.''
\end{quote}
\noindent --- Justice Harlan, dissenting, 403 U.S. at 752--753 (quoting Holmes, J.).

\end{document}
